% from ega2-1-fr.tex
\oldpage[II]{151}

\begin{remark}[7.4.12]
\label{II.7.4.12-fr}
Nous montrerons au chapitre V que la conclusion de
(\hyperref[II.7.4.10-fr]{7.4.10}) est encore valable sans supposer la courbe
normale (ni même réduite)~; nous montrerons aussi que pour qu'une courbe
algébrique (réduite ou non) soit affine, il faut et il suffit que ses
composantes irréductibles (réduites) ne soient pas complètes.
\end{remark}

\begin{corollary}[7.4.13]
\label{II.7.4.13-fr}
Soient $X$ une courbe normale irréductible de corps $R(X)=K$, $Y$ une courbe
intègre complète de corps $L=R(Y)$. Il y a une correspondance biunivoque
canonique entre les $k$-morphismes dominants $X\to Y$ et les
$k$-monomorphismes $L\to K$.
\end{corollary}

En vertu de (\hyperref[II.7.4.9-fr]{7.4.9}), les $k$-applications rationnelles
de $X$ dans $Y$ s'identifient aux $k$-morphismes $u:X\to Y$. Les morphismes
$u$ dominants étant caractérisés par le fait que $u(x)=y$ (en désignant par
$x$ et $y$ les points génériques respectifs de $X$ et de $Y$), le corollaire
résulte de ces remarques et de (\hyperref[I.7.1.13-fr]{I, 7.1.13}).

\begin{env}[7.4.14]
\label{II.7.4.14-fr}
On peut préciser le résultat de (\hyperref[II.7.4.13-fr]{7.4.13}) lorsqu'on
prend pour $Y$ la \emph{droite projective}
$\mathbf{P}_k^1=\operatorname{Proj}(k[T_0,T])$, $T_0$ et $T$ étant deux
indéterminées. C'est bien là un schéma intègre
(\hyperref[II.2.4.4-fr]{2.4.4}), et le schéma induit sur l'ouvert $D_+(T_0)$
de $Y$ est isomorphe à $\operatorname{Spec}(k[T])$
(\hyperref[II.2.3.6-fr]{2.3.6}), donc le point générique de $Y$ est l'idéal
$(0)$ de $k[T]$ et son corps de fonctions rationnelles $k(T)$, ce qui montre
que $Y$ est une courbe algébrique complète sur $k$. En outre, le seul idéal
premier gradué de $S=k[T_0,T]$ contenant $T_0$ et distinct de $S_+$ est
l'idéal principal $(T_0)$, donc le complémentaire de $D_+(T_0)$ dans
$Y=\mathbf{P}_k^1$ est réduit à \emph{un point fermé}, dit «~point à
l'infini~» que nous noterons $\infty$ (pour une étude générale des rapports
entre fibrés vectoriels et fibrés projectifs, voir
\hyperref[subsection:II.8.4-fr]{8.4}). Avec ces notations~:
\end{env}

\begin{corollary}[7.4.15]
\label{II.7.4.15-fr}
Soit $X$ une courbe normale irréductible de corps $R(X)=K$. Il existe une
application bijective canonique de $K$ sur l'ensemble des morphismes $u$ de
$X$ dans $\mathbf{P}_k^1$, distincts du morphisme constant de valeur
$\infty$. Pour que $u$ soit dominant, il faut et il suffit que l'élément
correspondant de $K$ soit transcendant sur $k$.
\end{corollary}

Cet énoncé résulte immédiatement de (\hyperref[II.7.4.9-fr]{7.4.9}) et du

\begin{lemma}[7.4.15.1]
\label{II.7.4.15.1-fr}
Soit $X$ un préschéma intègre sur $k$, et soit $K=R(X)$ son corps des fonctions
rationnelles. Il existe une application bijective canonique de l'ensemble
$K$ sur l'ensemble des applications rationnelles $u$ de $X$ dans
$\mathbf{P}_k^1$, distinctes du morphisme constant de valeur $\infty$. Pour
qu'une telle application rationnelle soit dominante, il faut et il suffit que
l'élément correspondant de $K$ soit transcendant sur $k$.
\end{lemma}

Tout d'abord, les applications rationnelles de $X$ dans $\mathbf{P}_k^1$
correspondent biunivoquement aux points de $\mathbf{P}_k^1$ à valeurs dans
l'extension $K$ de $k$ (\hyperref[I.7.1.12-fr]{I, 7.1.12}). Si un tel point
est localisé (\hyperref[I.3.4.5-fr]{I, 3.4.5}) au point générique de
$\mathbf{P}_k^1$, l'application rationnelle correspondante est évidemment
dominante. Dans le cas contraire, comme tout point de $\mathbf{P}_k^1$
distinct du point générique est fermé (\hyperref[II.7.4.3-fr]{7.4.3}),
l'image du domaine de définition $U$ de $u$ par l'unique morphisme
$U\to\mathbf{P}_k^1$ de la classe $u$ (\hyperref[I.7.2.2-fr]{I, 7.2.2}),
est réduite à un point fermé $y$ de $\mathbf{P}_k^1$, et ce morphisme (qui
n'est pas nécessairement partout défini dans $X$) n'est donc pas dominant~;
par abus de langage, on dit alors que l'application rationnelle $u$ est
«~constante, de valeur $y$~». Il reste à mettre en correspondance biunivoque
les points de $\mathbf{P}_k^1$ à valeurs dans $K$, de localité
(\hyperref[I.3.4.5-fr]{I, 3.4.5}) distincte de $\infty$, et les éléments de
$K$, et à vérifier que la localité d'un tel point est le point générique de
$\mathbf{P}_k^1$ si et seulement s'il correspond à un

